\chapter{Discussion and conclusion}
\label{chap:conclusion}

\section{What generalises}

The absolute numbers in \cref{chap:results} belong to one hardware
configuration and should not be quoted outside it. The \emph{shape} of
the curves is more robust: a fixed initialisation cost that caching
removes, a per-page cost that only incremental work removes, and a
crossover between the two that moves with document size exactly as
\cref{thm:crossover} predicts.

\section{Threats to validity}

\begin{description}
  \item[Synthetic documents.] Generated sources have a more uniform mix
        of content than real manuscripts. Documents dominated by large
        raster figures behave differently, because image handling is not
        page-proportional.
  \item[Single worker.] All measurements were taken on an unloaded
        worker. Contention between concurrent builds is left to future
        work \autocite{example2026conference}.
  \item[Median reporting.] Medians hide the tail that users actually
        complain about. \Cref{tab:percentiles} in \cref{chap:appendix-a}
        gives the 95th percentile for completeness.
\end{description}

\section{Conclusion}

Incremental recompilation is the only strategy of the three that keeps
build time roughly independent of document size, and it becomes the
better choice at around \num{60} pages -- well below the size of a
typical thesis. For collaborative systems, where the build runs hundreds
of times per session, that is the difference between a tool that keeps
up with the author and one that does not.

\begin{proposition}
Under the fitted parameters of \cref{tab:fit} and $\rho \le 0.1$, a
\num{300}-page document rebuilds incrementally in under \qty{2.5}{\second}.
\end{proposition}
